% Options for packages loaded elsewhere
\PassOptionsToPackage{unicode}{hyperref}
\PassOptionsToPackage{hyphens}{url}
\documentclass[
  11pt,
]{article}
\usepackage{xcolor}
\usepackage[margin=22mm]{geometry}
\usepackage{amsmath,amssymb}
\setcounter{secnumdepth}{-\maxdimen} % remove section numbering
\usepackage{iftex}
\ifPDFTeX
  \usepackage[T1]{fontenc}
  \usepackage[utf8]{inputenc}
  \usepackage{textcomp} % provide euro and other symbols
\else % if luatex or xetex
  \usepackage{unicode-math} % this also loads fontspec
  \defaultfontfeatures{Scale=MatchLowercase}
  \defaultfontfeatures[\rmfamily]{Ligatures=TeX,Scale=1}
\fi
\usepackage{lmodern}
\ifPDFTeX\else
  % xetex/luatex font selection
\fi
% Use upquote if available, for straight quotes in verbatim environments
\IfFileExists{upquote.sty}{\usepackage{upquote}}{}
\IfFileExists{microtype.sty}{% use microtype if available
  \usepackage[]{microtype}
  \UseMicrotypeSet[protrusion]{basicmath} % disable protrusion for tt fonts
}{}
\makeatletter
\@ifundefined{KOMAClassName}{% if non-KOMA class
  \IfFileExists{parskip.sty}{%
    \usepackage{parskip}
  }{% else
    \setlength{\parindent}{0pt}
    \setlength{\parskip}{6pt plus 2pt minus 1pt}}
}{% if KOMA class
  \KOMAoptions{parskip=half}}
\makeatother
\setlength{\emergencystretch}{3em} % prevent overfull lines
\providecommand{\tightlist}{%
  \setlength{\itemsep}{0pt}\setlength{\parskip}{0pt}}
\usepackage{bookmark}
\IfFileExists{xurl.sty}{\usepackage{xurl}}{} % add URL line breaks if available
\urlstyle{same}
\hypersetup{
  hidelinks,
  pdfcreator={LaTeX via pandoc}}

\author{}
\date{}

\begin{document}

\section{Inverting mixed symbols without commuting matrix
factors}\label{inverting-mixed-symbols-without-commuting-matrix-factors}

The inversion problem addressed here is pointwise in phase space. Its
output is a symbol with quantitative derivative estimates, so it can
subsequently enter an operator construction. No composition or
boundedness theorem for pseudodifferential operators is needed for the
inversion itself.

\href{euclidean-symbol-calculus.pdf}{From symbol estimates to operators
on every Sobolev scale} supplies the earlier symbol-space context. The
calculation below uses multivariable differentiation, the product rule,
matrix norms and smooth cutoffs; every mixed estimate needed here is
proved directly.

\subsection{1. The two frequency
weights}\label{the-two-frequency-weights}

Fix integers \(n,N\geq1\). Write

\[
x=(x_1,\ldots,x_n)\in\mathbb R^n,
\qquad
\xi=(\xi',\xi_n)\in\mathbb R^{n-1}\times\mathbb R.
\tag{MSI1}
\]

The matrix space is \(M_N(\mathbb C)\), equipped with the operator norm
induced by the Euclidean norm on \(\mathbb C^N\). Thus
\(\|UV\|\leq\|U\|\|V\|\), \(\|I_N\|=1\), and scalar multiplication has
its usual absolute-value norm. For \(n=1\), the vector \(\xi'\) has no
coordinates, its norm is zero, and every multiindex in its coordinates
has length zero.

For real numbers \(p,q\) and multiindices
\(\alpha,\beta\in\mathbb N_0^n\), put
\(\alpha'=(\alpha_1,\ldots,\alpha_{n-1})\) and

\[
W_{p,q;\alpha}(\xi)
=(1+|\xi|)^{p-\alpha_n}
 (1+|\xi'|)^{q-|\alpha'|}.
\tag{MSI2}
\]

The original weights \(1+|\xi|\) and \(1+|\xi'|\) are retained
throughout. They are positive numerical weights; differentiability of
these weights at the origin is neither assumed nor used.

A smooth function
\(a:\mathbb R_x^n\times\mathbb R_\xi^n\to M_N(\mathbb C)\) belongs to
\(S^{p,q}\) when every number

\[
\|a\|_{p,q;\beta,\alpha}
=\sup_{x,\xi}
\frac{\|\partial_x^\beta\partial_\xi^\alpha a(x,\xi)\|}
     {W_{p,q;\alpha}(\xi)}
\tag{MSI3}
\]

is finite. In particular the constants are uniform in all
\(x\in\mathbb R^n\); there is no implicit compact restriction on \(x\).
The convention \(D_j=-i\partial_j\), if used in an operator realization,
gives the same seminorms, because each derivative of fixed multiindex
differs by a scalar of modulus one. We use ordinary derivatives in the
identities below so that every sign is explicit.

The product of two smooth matrix symbols satisfies the exact Leibniz
identity

\[
\partial_x^\beta\partial_\xi^\alpha(uv)
=\sum_{\substack{\delta\leq\beta\\\varepsilon\leq\alpha}}
 \binom\beta\delta\binom\alpha\varepsilon
 (\partial_x^\delta\partial_\xi^\varepsilon u)
 (\partial_x^{\beta-\delta}\partial_\xi^{\alpha-\varepsilon}v).
\tag{MSI4}
\]

No factors on the right are interchanged. To prove this identity, apply
the ordinary one-coordinate product rule repeatedly; choosing which of
the \(\beta_j\) derivatives acts on the left factor gives
\(\binom{\beta_j}{\delta_j}\) choices, and the frequency derivatives
give the second collection of binomial factors. Multiplying these
independent counts gives (MSI4). Since

\[
W_{p,q;\varepsilon}
W_{r,t;\alpha-\varepsilon}
=W_{p+r,q+t;\alpha},
\tag{MSI5}
\]

submultiplicativity and (MSI4) prove that \(uv\in S^{p+r,q+t}\) whenever
\(u\in S^{p,q}\) and \(v\in S^{r,t}\). More precisely, its
\((\beta,\alpha)\) seminorm is bounded by the sum in (MSI4) with each
differentiated factor replaced by its corresponding seminorm. This
argument also applies on a subset of phase space whenever all
derivatives in the displayed estimates exist in a neighborhood of each
of its points.

\subsubsection{Exact identity map to bracket-weight
conventions}\label{exact-identity-map-to-bracket-weight-conventions}

Some statements use \(\langle\xi\rangle=(1+|\xi|^2)^{1/2}\) and
\(\langle\xi'\rangle=(1+|\xi'|^2)^{1/2}\). We retain the original
weights and prove their precise correspondence with these additional
weights. For every \(s\geq0\), \[
 1+s^2\leq(1+s)^2\leq2(1+s^2),
\] because the first difference is \(2s\geq0\) and the second difference
is \((s-1)^2\geq0\). Thus the two positive ratios \[
 \frac{1+|\xi|}{\langle\xi\rangle},\qquad
 \frac{1+|\xi'|}{\langle\xi'\rangle}
 \quad\hbox{belong to }[1,\sqrt2].
 \tag{MSI5a}
\] When \(n=1\), the second ratio is exactly one. For a real number
\(t\), put \(t_+=\max(t,0)\), \(t_-=\max(-t,0)\). Raising a number in
\([1,\sqrt2]\) to \(t\) gives a number in \([2^{-t_-/2},2^{t_+/2}]\);
for \(t<0\) the inequality reverses on taking the reciprocal, which
proves this assertion in that case as well.

Define the second, unchanged-derivative weight \[
 \widetilde W_{p,q;\alpha}(\xi)
 =\langle\xi\rangle^{p-\alpha_n}
  \langle\xi'\rangle^{q-|\alpha'|},
 \qquad u=p-\alpha_n,\quad v=q-|\alpha'|.
 \tag{MSI5b}
\] Multiplying the two ratio bounds gives, for these exact real
exponents, \[
 2^{-(u_-+v_-)/2}\widetilde W_{p,q;\alpha}
 \leq W_{p,q;\alpha}
 \leq2^{(u_++v_+)/2}\widetilde W_{p,q;\alpha}.
 \tag{MSI5c}
\] Let \(\widetilde{\|f\|}_{p,q;\beta,\alpha}\) denote (MSI3) with only
its denominator changed to \(\widetilde W\). Division by the positive
bounds in (MSI5c), followed by the same supremum over all \(x,\xi\),
proves \[
\begin{aligned}
 \|f\|_{p,q;\beta,\alpha}
 &\leq2^{(u_-+v_-)/2}\widetilde{\|f\|}_{p,q;\beta,\alpha},\\
 \widetilde{\|f\|}_{p,q;\beta,\alpha}
 &\leq2^{(u_++v_+)/2}\|f\|_{p,q;\beta,\alpha}.
\end{aligned}
 \tag{MSI5d}
\] Therefore the identity on the actual smooth matrix functions is a
continuous linear bijection between the two symbol spaces with their
families of seminorms. The two inequalities prove continuity in both
directions. It commutes with every derivative, preserves matrix
multiplication in its original order, and preserves the actual cutoff
inverse \(b\). These are literal identities on the functions, not
changes of phase-space coordinates or of the operator convention.

The ellipticity constants also have definite transformations. Write
\(c_R,C_R\) for scalar lower-bound and inverse-norm constants using
\(R=1+|\xi|\), and \(c_{\langle\cdot\rangle},C_{\langle\cdot\rangle}\)
for constants using \(\langle\xi\rangle\). The following choices are
valid on exactly the same set: \[
\begin{aligned}
 |a|\geq c_R R^m
   &\ \Longrightarrow\ c_{\langle\cdot\rangle}=c_R2^{-m_-/2},\\
 |a|\geq c_{\langle\cdot\rangle}\langle\xi\rangle^m
   &\ \Longrightarrow\ c_R=c_{\langle\cdot\rangle}2^{-m_+/2},\\
 \|a^{-1}\|\leq C_R R^{-m}
   &\ \Longrightarrow\ C_{\langle\cdot\rangle}=C_R2^{m_-/2},\\
 \|a^{-1}\|\leq C_{\langle\cdot\rangle}\langle\xi\rangle^{-m}
   &\ \Longrightarrow\ C_R=C_{\langle\cdot\rangle}2^{m_+/2}.
\end{aligned}
 \tag{MSI5e}
\] For example the first implication follows from
\(R^m\geq2^{-m_-/2}\langle\xi\rangle^m\); the second follows from
\(\langle\xi\rangle^m\geq2^{-m_+/2}R^m\). The last two apply the
corresponding upper bounds to exponent \(-m\). This proves all four
implications, including negative and zero orders. Together with (MSI5d),
it transfers every hypothesis and conclusion of the inverse theorem in
both directions with explicit constants, retaining both weight
presentations.

\subsection{2. An exact finite formula for inverse
derivatives}\label{an-exact-finite-formula-for-inverse-derivatives}

Use a combined multiindex \(\mu=(\beta,\alpha)\in\mathbb N_0^{2n}\), and
write \(\partial^\mu=\partial_x^\beta\partial_\xi^\alpha\). Inequalities
and sums of combined multiindices are componentwise;
\(\mu!=\prod_{j=1}^{2n}\mu_j!\), and \(0<|\nu|\) means that \(\nu\) is
nonzero.

Let \(a\) be smooth on an open set \(V\subset\mathbb R^{2n}\), with
every \(a(z)\) invertible. The inverse \(r=a^{-1}\) is smooth: the
cofactor formula expresses each entry of \(r\) as a polynomial in the
entries of \(a\) divided by \(\det a\), a nowhere-zero smooth function
on \(V\). For \(\mu\ne0\),

\[
\partial^\mu r
=\sum_{k=1}^{|\mu|}(-1)^k
 \sum_{\substack{\nu_1+\cdots+\nu_k=\mu\\|\nu_i|>0}}
 \frac{\mu!}{\nu_1!\cdots\nu_k!}
 r(\partial^{\nu_1}a)r\cdots
   (\partial^{\nu_k}a)r.
\tag{MSI6}
\]

The inner sum is over ordered lists. This is essential: in general, two
lists obtained by interchanging \(\nu_i\) and \(\nu_j\) produce
different matrix products.

Here is a finite algebraic proof of (MSI6). Fix \(z\in V\) and
\(L\geq|\mu|\). Work with polynomials in \(2n\) formal commuting
variables \(h=(h_1,\ldots,h_{2n})\), with matrix coefficients, and
identify two polynomials when their difference contains only monomials
of total degree greater than \(L\). For any smooth matrix function
\(f\), form the finite polynomial

\[
J_Lf(h)=\sum_{|\eta|\leq L}
 \frac{\partial^\eta f(z)}{\eta!}h^\eta.
\tag{MSI7}
\]

Formula (MSI4), coefficient by coefficient, proves
\(J_L(fg)=(J_Lf)(J_Lg)\) in this finite quotient algebra. Let
\(A=a(z)\), \(R=A^{-1}\), and \(H=J_La-A\). Every monomial of \(H\) has
positive degree. Consequently \((RH)^{L+1}=0\) in the quotient. Direct
multiplication, without commuting \(R\) and \(H\), now gives

\[
(A+H)^{-1}
=\sum_{k=0}^{L}(-RH)^kR.
\tag{MSI8}
\]

Indeed \(A+H=A(I_N+RH)\), and \((I_N+RH)\sum_{k=0}^{L}(-RH)^k=I_N\),
with the same identity in the reverse order. On the other hand,
\(J_La\,J_Lr=J_Lr\,J_La=I_N\), because \(ar=ra=I_N\) as actual smooth
functions. An inverse in any associative algebra is unique: if
\(XY=YX=I_N\) and \(XZ=ZX=I_N\), then \(Y=Y(XZ)=(YX)Z=Z\). Thus \(J_Lr\)
equals (MSI8). The coefficient of \(h^\mu\) in \((-RH)^kR\) is the
ordered product sum in (MSI6), divided by \(\mu!\). Comparing
coefficients proves the formula at the arbitrary point \(z\).

In particular, for a single coordinate \(z_j\) and for any two
coordinates \(z_j,z_l\),

\[
\partial_jr=-r(\partial_ja)r,
\qquad
\partial_l\partial_jr
=r(\partial_la)r(\partial_ja)r
 +r(\partial_ja)r(\partial_la)r
 -r(\partial_l\partial_ja)r.
\tag{MSI9}
\]

These identities hold for matrices with no symmetry, normality,
positivity, or commutation assumptions.

\subsection{3. A cutoff inverse in the full mixed
class}\label{a-cutoff-inverse-in-the-full-mixed-class}

\textbf{Theorem.} Let \(m\in\mathbb R\),
\(a\in S^{m,0}(\mathbb R^n\times\mathbb R^n;M_N(\mathbb C))\), and
\(\chi\in C_c^\infty(\mathbb R_\xi^n;\mathbb C)\). Write

\[
\theta(\xi)=1-\chi(\xi),\qquad
S=\operatorname{supp}\theta.
\tag{MSI10}
\]

Suppose that \(a(x,\xi)\) is invertible whenever \(x\in\mathbb R^n\) and
\(\xi\in S\), and that a constant \(C_0>0\) satisfies

\[
\|a(x,\xi)^{-1}\|
\leq C_0(1+|\xi|)^{-m}
\qquad (x\in\mathbb R^n,\ \xi\in S).
\tag{MSI11}
\]

Define the function on the whole phase space by

\[
b(x,\xi)=
\begin{cases}
\theta(\xi)a(x,\xi)^{-1},&\xi\in S,\\
0,&\xi\notin S.
\end{cases}
\tag{MSI12}
\]

Then \(b\in S^{-m,0}\), and the two pointwise identities

\[
a(x,\xi)b(x,\xi)=b(x,\xi)a(x,\xi)
=\theta(\xi)I_N
\tag{MSI13}
\]

hold on all of \(\mathbb R^{2n}\). In particular \(b=a^{-1}\) outside
\(\operatorname{supp}\chi\). The theorem imposes no invertibility
condition where \(\theta\) is identically zero in a neighborhood, and
imposes neither reality nor \(0\leq\chi\leq1\) on the cutoff.

For \(N=1\), the scalar hypothesis

\[
|a(x,\xi)|\geq c(1+|\xi|)^m
\quad (x\in\mathbb R^n,\ \xi\in S),\qquad c>0,
\tag{MSI14}
\]

is sufficient, with \(C_0=c^{-1}\). In that case (MSI12) is precisely
the smooth extension of \((1-\chi(\xi))/a(x,\xi)\) by zero on the open
region where \(1-\chi\) vanishes.

\textbf{Proof, including seminorm constants.} First we justify
smoothness of the globally defined function. If \(\xi_0\notin S\), then
\(\theta\) is identically zero on a neighborhood of \(\xi_0\), so \(b\)
is smooth near every \((x_0,\xi_0)\). If \(\xi_0\in S\), the invertible
matrix \(a(x_0,\xi_0)\) remains invertible on some open neighborhood of
\((x_0,\xi_0)\), because its continuous determinant is nonzero there
after shrinking the neighborhood. On that neighborhood the smooth
function \(\theta a^{-1}\) agrees with (MSI12): at points outside \(S\),
\(\theta=0\). Thus it is a local smooth representative of \(b\) at every
such point, including points of the boundary of \(S\). This proves
smoothness without imposing an additional neighborhood lower bound in
the hypotheses. Formula (MSI13) follows directly on \(S\); outside
\(S\), both sides are zero.

For every combined multiindex \(\nu=(\delta,\varepsilon)\), let

\[
A_\nu=\|a\|_{m,0;\delta,\varepsilon}.
\tag{MSI15}
\]

Define the nonnegative finite constants

\[
Q_0=C_0,
\qquad
Q_\mu=
\sum_{k=1}^{|\mu|}
 \sum_{\substack{\nu_1+\cdots+\nu_k=\mu\\|\nu_i|>0}}
 \frac{\mu!}{\nu_1!\cdots\nu_k!}
 C_0^{k+1}\prod_{i=1}^{k}A_{\nu_i}
\quad(\mu\ne0).
\tag{MSI16}
\]

At a point with \(\xi\in S\), the inverse derivatives are understood
through the local smooth inverse just constructed. Every term in (MSI6)
has \(k+1\) factors \(a^{-1}\) and \(k\) differentiated factors \(a\).
If \(\nu_i=(\beta^{(i)},\alpha^{(i)})\) and
\(\sum_i\nu_i=(\beta,\alpha)\), its total first-weight exponent is

\[
-(k+1)m+\sum_{i=1}^{k}(m-\alpha_n^{(i)})
=-m-\alpha_n;
\tag{MSI17}
\]

its second-weight exponent is \(-\sum_i|\alpha'^{(i)}|=-|\alpha'|\).
Therefore (MSI6), (MSI11), and the triangle inequality give the fully
uniform bounds

\[
\|\partial_x^\beta\partial_\xi^\alpha a^{-1}(x,\xi)\|
\leq Q_{(\beta,\alpha)}W_{-m,0;\alpha}(\xi)
\qquad (\xi\in S).
\tag{MSI18}
\]

For \(\alpha=\beta=0\), this is exactly (MSI11), so no derivative
induction lacks a starting case.

Choose a number \(R\geq1\) with
\(\operatorname{supp}\chi\subset\{|\xi|\leq R\}\); \(R=1\) is allowed
when \(\chi=0\). Set

\[
T_\delta=\sup_\xi|\partial_\xi^\delta\theta(\xi)|,
\qquad
L_\delta=
\begin{cases}1,&\delta=0,\\(1+R)^{|\delta|},&\delta\ne0.
\end{cases}
\tag{MSI19}
\]

All \(T_\delta\) are finite. For \(\delta\ne0\), the derivative
\(\partial^\delta\theta=-\partial^\delta\chi\) is supported in the ball
of radius \(R\). On that support,

\[
\frac{W_{-m,0;\alpha-\delta}(\xi)}
     {W_{-m,0;\alpha}(\xi)}
=(1+|\xi|)^{\delta_n}(1+|\xi'|)^{|\delta'|}
\leq (1+R)^{|\delta|}=L_\delta.
\tag{MSI20}
\]

For \(\delta=0\), this ratio equals one everywhere. The product rule for
\(b=\theta a^{-1}\), valid locally at every point of \(S\), now yields

\[
\|b\|_{-m,0;\beta,\alpha}
\leq
\sum_{\delta\leq\alpha}
 \binom\alpha\delta T_\delta L_\delta
 Q_{(\beta,\alpha-\delta)}.
\tag{MSI21}
\]

To see why this is also a global estimate, at every point outside \(S\)
all derivatives of \(b\) vanish on a neighborhood, and at every point of
\(S\) the preceding local calculation applies. This includes every
boundary point; no omitted limiting argument is needed. The right side
of (MSI21) is a finite number depending only on the explicitly listed
constants and derivative orders. Hence all seminorms required for
\(S^{-m,0}\) are finite. In the scalar case, (MSI14) gives (MSI11) by
taking reciprocals of the positive numerical lower bound. This proves
every assertion. \(\square\)

\subsection{4. Improvement for an individual frequency
derivative}\label{improvement-for-an-individual-frequency-derivative}

\textbf{Theorem.} Retain all hypotheses of Section 3. Fix one index
\(j\in\{1,\ldots,n\}\). If

\[
a_j:=\partial_{\xi_j}a\in S^{m-1,0},
\tag{MSI22}
\]

then

\[
\partial_{\xi_j}b\in S^{-m-1,0}.
\tag{MSI23}
\]

There is no hypothesis here on the other first frequency derivatives
beyond \(a\in S^{m,0}\). Consequently the same conclusion holds for any
collection of indices for which (MSI22) holds, including all indices
when the stronger assumption holds for all of them.

\textbf{Proof.} At every point with \(\xi\in S\), the exact matrix
identity (MSI9) gives

\[
\partial_{\xi_j}b
=-\theta a^{-1}a_j a^{-1}
  +(\partial_{\xi_j}\theta)a^{-1}.
\tag{MSI24}
\]

The order of the two inverses and \(a_j\) is fixed. Each term has a
smooth global extension by zero outside \(S\). Indeed near a point of
\(S\) it is the displayed product of smooth functions and the locally
defined smooth inverse. At a point outside \(S\), \(\theta\), all its
derivatives, and the extended terms vanish on a neighborhood. On an
overlap these local definitions agree. For the term involving
\(\partial_{\xi_j}\theta\), observe that every derivative of \(\theta\)
vanishes outside \(S\), because \(\theta\) is locally zero there. Thus
the same gluing argument applies to that term.

We give explicit bounds for every derivative in (MSI24). Let

\[
H^{(j)}_\kappa
=\|a_j\|_{m-1,0;\kappa_x,\kappa_\xi}
\quad
\text{for }\kappa=(\kappa_x,\kappa_\xi)\in\mathbb N_0^{2n}.
\tag{MSI25}
\]

When \(\mu=(\beta,\alpha)\), a term in the repeated product rule for
\(\theta a^{-1}a_j a^{-1}\) is indexed by

\[
(0,\delta)+\nu+\kappa+\lambda=\mu,
\qquad \delta\in\mathbb N_0^n,
\quad \nu,\kappa,\lambda\in\mathbb N_0^{2n}.
\tag{MSI26}
\]

The differentiated three matrix factors have combined first exponent
\(-m+(m-1)-m-(\alpha_n-\delta_n)=-m-1-\alpha_n+\delta_n\), and combined
second exponent \(-|\alpha'|+|\delta'|\). If \(\delta=0\), these are
exactly the desired exponents. If \(\delta\ne0\), the cutoff derivative
restricts the term to \(|\xi|\leq R\), and the weight correction is
bounded by \(L_\delta\), exactly as in (MSI20). Thus the
\((\beta,\alpha)\) seminorm of the first term of (MSI24), in
\(S^{-m-1,0}\), is at most

\[
U_\mu^{(j)}=
\sum_{(0,\delta)+\nu+\kappa+\lambda=\mu}
 \frac{\mu!}{\delta!\nu!\kappa!\lambda!}
 T_\delta L_\delta Q_\nu H^{(j)}_\kappa Q_\lambda.
\tag{MSI27}
\]

Here \((0,\delta)!=\delta!\), which accounts for the denominator. The
finite multinomial count follows from the same derivative-choice
argument that proved (MSI4), now with four factors.

For the second term of (MSI24), apply the product rule with
\(\delta\leq\alpha\) derivatives hitting \(\partial_{\xi_j}\theta\). All
such cutoff factors are supported in \(|\xi|\leq R\), even for
\(\delta=0\), because their total derivative order \(|\delta|+1\) is
positive. Their weight comparison is

\[
\frac{W_{-m,0;\alpha-\delta}(\xi)}
     {W_{-m-1,0;\alpha}(\xi)}
=(1+|\xi|)^{1+\delta_n}
 (1+|\xi'|)^{|\delta'|}
\leq(1+R)^{1+|\delta|}.
\tag{MSI28}
\]

If \(e_j\in\mathbb N_0^n\) is the \(j\)-th frequency multiindex, the
second term therefore has seminorm at most

\[
V_\mu^{(j)}=
\sum_{\delta\leq\alpha}
 \binom\alpha\delta T_{\delta+e_j}
 (1+R)^{1+|\delta|}Q_{(\beta,\alpha-\delta)}.
\tag{MSI29}
\]

Combining (MSI24), (MSI27), and (MSI29) gives

\[
\|\partial_{\xi_j}b\|_{-m-1,0;\beta,\alpha}
\leq U_{(\beta,\alpha)}^{(j)}+V_{(\beta,\alpha)}^{(j)}<\infty.
\tag{MSI30}
\]

The local extension argument again makes these global bounds. This
proves (MSI23) for the chosen \(j\), without invoking any condition on
another first derivative. \(\square\)

There is an immediate distinction between the normal and tangential
indices. The normal case \(j=n\) of (MSI22) follows already from
\(a\in S^{m,0}\), since

\[
W_{m,0;\alpha+e_n}=W_{m-1,0;\alpha}.
\tag{MSI31}
\]

For \(j<n\), the original class instead yields
\(\partial_{\xi_j}a\in S^{m,-1}\). The first-frequency weight and
second-frequency weight have different behavior when
\(|\xi_n|\to\infty\) while \(\xi'\) is fixed. The following explicit
example proves that the tangential improvement cannot be inferred from
the original class alone.

\subsection{5. Two examples with different
purposes}\label{two-examples-with-different-purposes}

\textbf{A singular matrix in the discarded region.} Let \(N=2\),
\(n\geq1\), \(m=2\), and

\[
J=\begin{pmatrix}0&1\\0&0\end{pmatrix},
\qquad
a(x,\xi)=|\xi|^2I_2+J.
\tag{MSI32}
\]

This symbol is independent of \(x\), and \(J^2=0\), \(\|J\|=1\). It
belongs to \(S^{2,0}\), as may be checked at every derivative order
directly. At order zero, \(\|a\|\leq|\xi|^2+1\leq(1+|\xi|)^2\). At a
first frequency derivative, \(\partial_{\xi_j}a=2\xi_jI_2\). For
\(j=n\), this is bounded by \(2(1+|\xi|)\), the required normal weight.
For \(j<n\), \(|\xi_j|\leq|\xi'|\), so

\[
2|\xi_j|(1+|\xi'|)
\leq2|\xi'|(1+|\xi'|)
\leq2(1+|\xi|)^2.
\tag{MSI33}
\]

At order two, the only nonzero derivatives are
\(\partial_{\xi_j}^2a=2I_2\). For \(j=n\), the prescribed weight is one.
For \(j<n\), it is \((1+|\xi|)^2(1+|\xi'|)^{-2}\geq1\). All mixed second
derivatives, all derivatives of order greater than two, and every
positive-order \(x\) derivative vanish. These observations prove all
required seminorm bounds.

Every \(a_j=2\xi_jI_2\) also lies in \(S^{1,0}\). At order zero its norm
is at most \(2(1+|\xi|)\). Its only possibly nonzero positive-order
derivative is \(\partial_{\xi_j}a_j=2I_2\); the required weight is one
if \(j=n\), and \((1+|\xi|)/(1+|\xi'|)\geq1\) if \(j<n\). All other
derivatives vanish.

Choose a smooth cutoff \(\chi(\xi)\) equal to one when \(|\xi|\leq1\),
equal to zero when \(|\xi|\geq2\), and taking values in \([0,1]\). Such
a cutoff can be made explicitly. Define \(h(t)=e^{-1/t}\) for \(t>0\)
and \(h(t)=0\) for \(t\leq0\); repeated differentiation for \(t>0\)
gives a polynomial in \(t^{-1}\) times \(e^{-1/t}\), which tends to zero
as \(t\downarrow0\), at every derivative order. Thus \(h\) is smooth.
Set

\[
\chi(\xi)
=\frac{h(4-|\xi|^2)}
       {h(4-|\xi|^2)+h(|\xi|^2-1)}.
\tag{MSI34}
\]

The denominator is strictly positive for every \(\xi\): if
\(|\xi|^2\leq1\), the first term is positive; if \(|\xi|^2\geq4\), the
second is positive; in between both are positive. This proves smoothness
and all stated cutoff properties, including compact support. Every
\(\xi\in\operatorname{supp}(1-\chi)\) has \(|\xi|\geq1\). For such
\(\xi\), direct multiplication using \(J^2=0\) gives

\[
a(x,\xi)^{-1}=|\xi|^{-2}I_2-|\xi|^{-4}J,
\qquad
\|a(x,\xi)^{-1}\|
\leq2|\xi|^{-2}\leq8(1+|\xi|)^{-2}.
\tag{MSI35}
\]

Thus (MSI11) holds with \(C_0=8\). The resulting \(b\) belongs to
\(S^{-2,0}\), and every frequency derivative belongs to \(S^{-3,0}\).
Although \(a(x,0)=J\) is singular, \(b=0\) on \(|\xi|\leq1\), and the
globally smooth inverse construction remains valid. This verifies
concretely why global invertibility of \(a\) is unnecessary.

\textbf{A tangential derivative that does not improve.} Let \(n\geq2\),
and choose a real-valued \(\varphi\in C_c^\infty(\mathbb R^{n-1})\) with
\(|\varphi|\leq1/2\) and \(\partial_{\xi_1}\varphi(\eta)\ne0\) at some
point \(\eta\). Such a function is obtained by multiplying the
coordinate \(\xi_1\) by a smooth cutoff equal to one near the origin and
then multiplying by a sufficiently small positive constant; the cutoff
is constructed by the same one-variable function \(h\) as in (MSI34).
Put

\[
a(x,\xi)=2+\varphi(\xi'),\qquad \chi=0,
\qquad b(x,\xi)=\frac1{2+\varphi(\xi')}.
\tag{MSI36}
\]

The symbol \(a\) is scalar and at least \(3/2\). To verify
\(a\in S^{0,0}\), if \(\alpha_n>0\) or \(\beta\ne0\), the required
derivative is zero. If \(\beta=0\), \(\alpha_n=0\), and \(|\alpha'|>0\),
the derivative \(\partial_{\xi'}^{\alpha'}\varphi\) has compact support
in \(\xi'\). If that support lies in \(|\xi'|\leq R_\varphi\), the
required seminorm is at most
\(\|\partial^{\alpha'}\varphi\|_\infty(1+R_\varphi)^{|\alpha'|}\). The
order-zero seminorm is at most \(5/2\). Thus every required estimate is
established. The inverse bound (MSI11) holds for \(m=0\) with
\(C_0=2/3\), and Section 3 yields \(b\in S^{0,0}\).

At \(\xi'=\eta\), however,

\[
\partial_{\xi_1}b(x,\eta,\xi_n)
=-\frac{\partial_{\xi_1}\varphi(\eta)}
        {(2+\varphi(\eta))^2}\ne0,
\tag{MSI37}
\]

independently of \(\xi_n\). The order-zero estimate for membership in
\(S^{-1,0}\) would bound the absolute value of this nonzero constant by
\(C(1+\sqrt{|\eta|^2+\xi_n^2})^{-1}\) for all \(\xi_n\). The right side
tends to zero as \(|\xi_n|\to\infty\), which is impossible. Therefore
\(\partial_{\xi_1}b\notin S^{-1,0}\). The same argument shows
\(\partial_{\xi_1}a\notin S^{-1,0}\). This example preserves the full
original mixed class while proving that the extra tangential assumption
in (MSI22) has mathematical content.

\subsection{6. Solved exercise: detect an invalid
commutation}\label{solved-exercise-detect-an-invalid-commutation}

\textbf{Exercise.} For \(n\geq1\), take a matrix symbol independent of
\(\xi\) and of \(x_2,\ldots,x_n\), defined by

\[
a(x,\xi)=
\begin{pmatrix}1&\sin x_1\\\cos x_1&1\end{pmatrix},
\qquad \chi=0.
\tag{MSI38}
\]

Prove that Section 3 applies with \(m=0\). Compute \(b\),
\(\partial_{x_1}b\), and \(\partial_{x_1}^2b\) at \(x_1=0\). Determine
whether replacing the inverse derivative by \(- (\partial_{x_1}a)b^2\)
gives the correct first derivative there. Finally, explain the
frequency-derivative conclusion for this example.

\textbf{Solution.} Every \(x_1\) derivative of every entry of \(a\) is
bounded, and all positive-order frequency derivatives vanish. These
facts verify every seminorm of \(S^{0,0}\): when \(\alpha=0\), the
denominator in (MSI3) is one; when \(\alpha\ne0\), the numerator
vanishes. The determinant is

\[
d(x_1)=1-\sin x_1\cos x_1\geq\tfrac12,
\qquad
b(x,\xi)=\frac1{d(x_1)}
\begin{pmatrix}1&-\sin x_1\\-\cos x_1&1\end{pmatrix}.
\tag{MSI39}
\]

The bound on the determinant follows from
\(2|\sin x_1\cos x_1|\leq\sin^2x_1+\cos^2x_1=1\). The numerator matrix
has squared Frobenius norm \(2+\sin^2x_1+\cos^2x_1=3\), so its Euclidean
operator norm is at most \(\sqrt3\): for any vector \(v\), the scalar
Cauchy--Schwarz inequality applied to each row gives
\(\|Mv\|^2\leq(\sum_{k,l}|M_{kl}|^2)\|v\|^2\). Consequently
\(\|b\|\leq2\sqrt3\), and (MSI11) holds with \(C_0=2\sqrt3\). In
particular the theorem supplies all inverse-symbol seminorms without
assuming them in advance.

At \(x_1=0\), write

\[
A=\begin{pmatrix}1&0\\1&1\end{pmatrix},
\quad R=A^{-1}=\begin{pmatrix}1&0\\-1&1\end{pmatrix},
\quad A_1=\begin{pmatrix}0&1\\0&0\end{pmatrix},
\quad A_2=\begin{pmatrix}0&0\\-1&0\end{pmatrix}.
\tag{MSI40}
\]

Formula (MSI9), with the products performed in their written order,
gives

\[
\left.\partial_{x_1}b\right|_{x_1=0}
=-RA_1R
=\begin{pmatrix}1&-1\\-1&1\end{pmatrix},
\tag{MSI41}
\]

and

\[
\left.\partial_{x_1}^2b\right|_{x_1=0}
=2RA_1RA_1R-RA_2R
=\begin{pmatrix}2&-2\\-1&2\end{pmatrix}.
\tag{MSI42}
\]

For a direct check of the first multiplication,
\(RA_1R=\bigl(\begin{smallmatrix}-1&1\\1&-1\end{smallmatrix}\bigr)\).
For the second,
\(RA_1RA_1R=\bigl(\begin{smallmatrix}1&-1\\-1&1\end{smallmatrix}\bigr)\)
and \(RA_2R=\bigl(\begin{smallmatrix}0&0\\-1&0\end{smallmatrix}\bigr)\);
these two explicit products yield (MSI42).

By contrast,

\[
-A_1R^2
=\begin{pmatrix}2&-1\\0&0\end{pmatrix}
\ne
\begin{pmatrix}1&-1\\-1&1\end{pmatrix}.
\tag{MSI43}
\]

The proposed commutation therefore fails even for a uniformly
invertible, smooth, globally bounded symbol. Every positive-order
frequency derivative of \(a\) and \(b\) is zero. The symbols themselves
have bounded position derivatives of every order, so both belong to
\(S^{0,0}\). Indeed their entries are rational functions of sine and
cosine with denominators bounded away from zero; repeated position
differentiation therefore remains bounded. The nonzero symbols are not
asserted to belong to every mixed order class. In particular (MSI22) and
(MSI23) hold for each \(j\in\{1,\ldots,n\}\) with \(m=0\). This last
conclusion comes from actual vanishing, whereas the noncommuting \(x_1\)
derivatives above remain nonzero. \(\square\)

\subsection{References}\label{references}

\end{document}
